\section{Finite states and the information required at the endpoint}
\label{sec:model-endpoint}

The block estimates will bound information in a state after an independent
projection is known. After specifying the learner model, we show why an
accurate state must still carry order $d\log(1/\epsilon)$ information in
that experiment. We also record
the reductions for randomized rules, rare states, and bounded stopping,
so that later applications need only specify their block dimensions and
their way of charging the clock.

\subsection{Finite-state learners}

\begin{definition}[Finite-state noiseless regression]\label{def:streaming-model}
An unknown unit vector $s\in S^{d-1}$ is fixed before sampling. At index
$j\ge1$ the learner receives
\[
 X_j\sim N(0,I_d),\qquad Y_j=\langle X_j,s\rangle,
\]
where the $X_j$ are independent. It cannot choose a row or revisit an earlier
one. At each index its persistent state has at most $2^M$ values. Its
transition and stopping rule may be any measurable randomized kernel of
the current state and the entire current pair $(X_j,Y_j)$. Computation
within a transition is unrestricted, and every piece of information retained
for later indices must be in the finite state.

The learner stops at an index $\tau\in\{0,\ldots,T\}$, where the integer
$T$ is a prescribed deterministic horizon. Its unit-vector output is drawn
using only the terminal state, the rule at index $\tau$, and fresh randomness.
In particular, the last observation can affect the output only through the
terminal state. Rules may depend on $d$, the target accuracy, and the index,
and may be selected by a shared random seed. The initial state and that
seed are jointly independent of the signal and the entire sample sequence.
Fresh transition and output random variables may also be used; each is
independent of that pair, the signal, the sample sequence, and all earlier
fresh random variables before it is used. Memory is measured in bits,
not real registers.

The construction from the signal, pre-generated sample rows, shared seed,
initialization, and fresh transition and output randomness through the
state path, stopping index, and output must form a jointly measurable
probability experiment. Per-seed measurability of the sample kernels alone
is not a definition of the mixture over seeds.
\end{definition}

For each fixed seed, the transition and stopping coordinates may be
measurable in the completion of the Borel sigma field under
$\lambda(dx,dy)=\gamma_{1,d}(dx)\,dy$ on $\R^d\times\R$, where
$\gamma_{1,d}$ is standard Gaussian probability on $\R^d$.
Lemma~\ref{lem:borel-rules} supplies Borel versions for the fixed-seed
uniform-prior experiment. The original randomized experiment must still
be jointly measurable as stated in the definition.

\subsection{Spherical caps and residual fibers}

\begin{lemma}[Spherical balls]\label{lem:spherical-ball}
Let $\sigma_q$ be uniform probability on $S^{q-1}$, where $q\ge2$.
For every $z\in S^{q-1}$ and $r>0$,
\begin{equation}\label{eq:spherical-ball}
 \sigma_q\{u:\norm{u-z}\le r\}\le r^{q-1}.
\end{equation}
In particular an angular cap of radius $0<\epsilon\le\pi$ has mass at most
$\epsilon^{q-1}$.
\end{lemma}

\begin{proof}
For $r\le1$, rotate $z$ to the north pole. A point in the chordal ball has
last coordinate at least $1-r^2/2\ge1/2$. As a graph over its first $q-1$
coordinates, the surface has area element at most twice Lebesgue measure,
and its projection lies in the $(q-1)$-ball of radius $r$. The whole sphere
has area at least twice the volume of the unit $(q-1)$-ball, by projecting
both hemispheres. This proves \eqref{eq:spherical-ball}. For $r>1$, total
mass one suffices. Angular distance at most $\epsilon$ implies chordal
distance $2\sin(\epsilon/2)\le\epsilon$, proving the last assertion.
\end{proof}

\begin{lemma}[Residual sphere and radius bounds]\label{lem:residual-sphere}
Let $S\sim\sigma_d$ and let $G$ be independent of $S$, with $b$ rows,
where $0\le b\le d-2$. For $b\ge1$, let $G$ be either a standard Gaussian
matrix or a uniform Haar frame of $b$ orthonormal rows in $\R^d$. For $b=0$,
interpret $G$ and $GS$ as empty. Conditional on $(G,GS)$, the signal is uniform on an affine sphere
parallel to $\ker G$, with ambient kernel dimension $d-b$ and radius
\[
 R=\norm{P_{\ker G}S}.
\]
For $b=0$, this is the original unit sphere and $R=1$ deterministically.
For $1\le b\le d-2$, the squared radius has distribution
\begin{equation}\label{eq:residual-beta}
 R^2\sim\operatorname{Beta}\bigl((d-b)/2,b/2\bigr).
\end{equation}
For $0<r_0<1$ one has
\begin{equation}\label{eq:residual-markov}
 \Prob\{R<r_0\}\le \frac{b}{d(1-r_0^2)}.
\end{equation}
When $0\le b\le d/2$, two useful bounds at radius $1/2$ are
\begin{equation}\label{eq:residual-chebyshev}
 \Prob\{R<1/2\}\le\frac{8}{d+2},
 \qquad
 \Prob\{R<1/2\}\le\frac{90}{d}.
\end{equation}
The second inequality also follows directly from two Gaussian norm
estimates, without using the beta variance.
\end{lemma}

\begin{proof}
The zero-row case is immediate: conditioning on empty data changes nothing,
and the residual radius is one. It does not involve a beta law with a zero
shape parameter.

Suppose $b\ge1$. The matrix is full row rank almost surely in the Gaussian
case and by definition in the Haar case. For fixed $G$,
write $S=Z/\norm Z$ with $Z$ standard Gaussian and split $Z$ into the row
space of $G$ and its kernel. The squared lengths are independent
chi-squares with $b$ and $d-b$ degrees of freedom. Their normalized ratio
gives \eqref{eq:residual-beta}; the direction in the kernel is independent
and uniform. Given the projection, its row-space component is fixed, so
the remaining signal is uniform on the asserted affine sphere. Its radius
has
\[
 \E R^2=\frac{d-b}{d},\qquad
 \operatorname{Var}(R^2)=\frac{2b(d-b)}{d^2(d+2)}.
\]
These formulas also hold for the deterministic zero-row radius. Markov's
inequality applied to $1-R^2$, whose mean is $b/d$, proves
\eqref{eq:residual-markov}. If $b\le d/2$, the mean of $R^2$ is at least
$1/2$ and its variance is at most $1/(2(d+2))$. Chebyshev's inequality
at distance $1/4$ from the mean gives the first bound in
\eqref{eq:residual-chebyshev}.

For the second bound, still with $b\le d/2$, let $Z_{\ker}$ be the kernel
component of $Z$. Its squared norm has mean $d-b\ge d/2$ and variance at
most $2d$, while $\norm Z^2$ has mean $d$ and variance $2d$. Chebyshev gives
\[
 \Prob\{\norm{Z_{\ker}}^2<d/3\}\le\frac{72}{d},\qquad
 \Prob\{\norm Z^2>4d/3\}\le\frac{18}{d}.
\]
Outside these events $R^2\ge1/4$. Their union gives $90/d$. For $b=0$
the probability is zero, as already established.
\end{proof}

The next estimate tests success against a law in which the signal and
output are conditionally independent. This is the binary-relative-entropy
event method; see Bassily et al.~\cite[Appendix~A.1, Lemmas~14--15]{BassilyEtAl2018}
for finite-space information bounds. Here the reference signal lies on
a residual sphere.

\begin{lemma}[Information on a residual sphere]\label{lem:residual-information}
Let $S\sim\sigma_d$, let $V$ be a random variable with values in a standard
Borel space, jointly distributed with $S$, and let $G$ be independent of
$(S,V)$, with $0\le b\le d-2$ rows. For $b\ge1$, $G$ may be either a standard
Gaussian matrix or a uniform Haar frame of $b$ orthonormal rows; at $b=0$
it is empty. Let $\widehat S\in S^{d-1}$
be generated from $V$ by a kernel conditionally independent of $(S,G)$.
Write
\[
 p=\Prob\{\arccos\langle\widehat S,S\rangle\le\epsilon\},
 \qquad R=\norm{P_{\ker G}S}.
\]
For $2\epsilon<r_0\le1$, if $\alpha:=p-\Prob\{R<r_0\}>0$, then
\begin{equation}\label{eq:residual-information}
 I(S;V\mid G,GS)\ge
 \alpha(d-b-1)\log\frac{r_0}{2\epsilon}-\log2.
\end{equation}
\end{lemma}

\begin{proof}
Put $Z=(G,GS)$. Compare the true law of $(S,V,Z)$ with
$P_ZP_{S\mid Z}P_{V\mid Z}$, and in both laws generate $\widehat S$ using
the same kernel from $V$. Adding that kernel preserves the divergence,
which is $I(S;V\mid Z)$.

Fix side data with residual radius $R\ge r_0$ and fix an output. If the
set of successful points on the fiber is nonempty, choose one such point.
Angular success implies Euclidean distance at most $\epsilon$ from the
output, so any two successful fiber points are at Euclidean distance at
most $2\epsilon$. After translating and rescaling the fiber to its unit
sphere, the successful set therefore lies in a chordal ball of radius
$2\epsilon/R\le2\epsilon/r_0<1$ centered on that sphere. By
Lemma~\ref{lem:spherical-ball}, its conditional mass is at most
$(2\epsilon/r_0)^{d-b-1}$. Under the comparison law, the output and signal
are conditionally independent given $Z$, so the event of success and
$R\ge r_0$ has probability at most this quantity. Under the true law it has
probability at least $\alpha$.

For Bernoulli probabilities $a,c$, relative entropy obeys
\[
 \KL(\operatorname{Bern}(a)\Vert\operatorname{Bern}(c))
 \ge a\log(1/c)-\log2;
\]
this follows by dropping the nonnegative term
$(1-a)\log(1/(1-c))$ and using binary entropy at most $\log2$.
Apply data processing to the event just described. If the comparison
probability is zero the conclusion is immediate; otherwise its upper bound
and the true lower bound give \eqref{eq:residual-information}.
\end{proof}

We record three numerical forms used below. The first uses an angular
cap bound: a chordal radius $4\epsilon$ corresponds to angular radius
$2\arcsin(2\epsilon)\le5\epsilon$ for $\epsilon\le1/10$.

\begin{corollary}[Three endpoint substitutions]\label{cor:replica-endpoints}
Assume $0<\epsilon\le1/10$ in Lemma~\ref{lem:residual-information}.
For sufficiently large $d$ the following bounds hold.
\begin{enumerate}[label=(\roman*)]
\item If $G$ is Gaussian, $p\ge3/5$, and $0\le b\le d/2$, the beta-variance bound (with
the separate deterministic case at $b=0$) and the angular cap just described
give
\begin{equation}\label{eq:endpoint-gaussian}
 I(S;V\mid G,GS)\ge
 \tfrac12(d-b-1)\log\frac1{5\epsilon}-\log2.
\end{equation}
\item If $G$ is a Haar frame, $p\ge2/3$, and $b=\lfloor d/4\rfloor$, Markov's bound
$\Prob\{R<1/2\}\le4b/(3d)\le1/3$ gives
\begin{equation}\label{eq:endpoint-haar}
 I(S;V\mid G,GS)\ge
 \tfrac13(d-b-1)\log\frac1{4\epsilon}-\log2.
\end{equation}
\item If $G$ is Gaussian, $p\ge3/5$, and $b=\lfloor d/2\rfloor$, the direct Gaussian bound
$\Prob\{R<1/2\}\le90/d$ gives
\begin{equation}\label{eq:endpoint-bins}
 I(S;V\mid G,GS)\ge
 \tfrac12(d-b-1)\log\frac1{4\epsilon}-\log2.
\end{equation}
\end{enumerate}
Each right side is at least $c\,d\log(1/\epsilon)$ for an absolute $c>0$
and all sufficiently large $d$.
\end{corollary}

\begin{proof}
Use $r_0=1/2$. In (i), the failure probability is at most $8/(d+2)$, so
the true event has mass at least $1/2$ for large $d$. The angular cap bound
is $(5\epsilon)^{d-b-1}$, which gives the displayed version of the binary
entropy calculation. In (ii), Markov's estimate leaves event mass at least
$1/3$, and the chordal cap is $(4\epsilon)^{d-b-1}$. In (iii), the stated
$90/d$ estimate leaves mass at least $1/2$. The last assertion follows from
$d-b-1\ge d/2-1$ and
\[
 \log\frac1{c_0\epsilon}\ge
 \left(1-\frac{\log c_0}{\log10}\right)\log\frac1\epsilon
 \qquad(c_0=4\text{ or }5).
\]
The positive multiple of $d\log(1/\epsilon)$ absorbs $\log2$ eventually.
\end{proof}

\subsection{Borel rules, randomness, and bounded stopping}

\begin{lemma}[Borel versions for the averaged experiment]\label{lem:borel-rules}
Fix $d\ge2$, a finite horizon $T$, and finite state sets. For fixed
data-independent rule parameters, suppose each transition and stopping
coordinate is measurable in the completion of the Borel sigma field on
$\R^d\times\R$ under
$\lambda(dx,dy)=\gamma_{1,d}(dx)\,dy$, where $\gamma_{1,d}$ is standard
Gaussian probability on $\R^d$. The output law depends only on the
terminal state and stopping index. There are Borel transition and stopping
kernels, with the same state sets and output laws, having the same joint law
of the signal, state path, stopping index, and output when $S\sim\sigma_d$.
\end{lemma}

\begin{proof}
For each state and index combine the next-state and stopping alternatives
into a finite destination set. Choose a Borel version of every coordinate
of this finite probability vector. Such versions exist for completed-
measurable real functions: approximate by simple functions and replace
their countably many level sets by Borel sets modulo null sets. On the Borel
null set where the resulting vector is negative in a coordinate or fails
to sum to one, replace it by a fixed point mass. The finite union over all
states and indices is contained in one Borel $\lambda$-null set $N$.

For a nonzero fixed $x$, the scalar $\langle x,S\rangle$ has a Lebesgue
density when $S\sim\sigma_d$ and $d\ge2$. Rotation and the one-coordinate
spherical law show this directly. The event $x=0$ is Gaussian-null.
Consequently each pre-generated sample pair
$(X_j,\langle X_j,S\rangle)$ has law absolutely continuous with respect
to $\lambda$, and it avoids $N$ almost surely. This remains true on the
event that any particular state is reached, because that event intersected
with the pair lying in $N$ is contained in the marginal null event.
The finite union over sample indices is still null. Coupling each original
finite kernel and its Borel version by the same fresh uniform random
variable now makes their state paths and stopping decisions agree by
induction. Their output laws agree at the identical terminal pair.
\end{proof}

This lemma preserves the experiment averaged over the uniform signal; it
does not claim pointwise preservation at every fixed signal. We may therefore
apply it after passing to the uniform-prior hypothesis of
Theorem~\ref{thm:streaming}.
When a shared seed is present, apply the replacement separately in its
fixed-seed experiments. The original jointly measurable law supplies
well-defined conditional success probabilities and their average. No
jointly measurable selection of the replacement Borel versions is needed.

\begin{lemma}[Fixing randomness and padding the clock]\label{lem:padding}
Consider a learner in Definition~\ref{def:streaming-model} under the uniform
prior. If its average success is at least $p$, some fixed shared-seed
conditional experiment has success at least $p$. Its initialization remains
independent of the signal and samples, and its fixed-seed Borel replacement
has the same success. In this Borel experiment the initialization and all
remaining randomized rules can be represented by data-independent variables
chosen before sampling, followed by deterministic state-based rules. For
every $p'<p$, some realization of those variables gives deterministic rules
with average success at least $p'$.

Padding a learner stopped by $T$ so that its halted states are absorbing,
while retaining the stopping index and postponing the output, uses at most
$(T+2)2^M$ states at every layer. Hence every padded state has entropy at
most
\begin{equation}\label{eq:padded-entropy}
 h_T=M\log2+\log(T+2).
\end{equation}
Extra ignored samples may be generated to complete a last block.
\end{lemma}

\begin{proof}
The original jointly measurable law gives a measurable conditional success
function of the shared seed whose average is the original success. Work on
a common full-measure set of seed values where the conditional law is the
stated fixed-seed experiment and the conditional initialization remains
independent of the signal and rows. The latter property follows from the
joint independence of the initial state and shared seed from the signal and
all rows. Some seed in this set has conditional success at least $p$:
otherwise the strictly positive difference between $p$ and that conditional
success would have a positive integral. Apply
Lemma~\ref{lem:borel-rules} to this one fixed-seed experiment. Its
uniform-prior success is unchanged. This step does not require a jointly
measurable choice of Borel replacements across seeds.

In the resulting fixed-seed Borel experiment, for a finite destination
kernel use an independent uniform variable and
the cumulative coordinate probabilities to choose the destination. There
are only finitely many state/index rules; unused choices can be generated
in advance as well. Similarly pre-generate, for each terminal state/index,
an output with its prescribed distribution, independently of the data.
Together with the conditional initial state, these variables are jointly
independent of the signal and samples. Conditional on their realization,
the initialization and rules are deterministic and state based. The finite
Borel construction is jointly measurable, and averaging its success gives
that of the selected fixed-seed experiment. Some realization therefore has
success at least $p'<p$.

At a layer the active state set has at most $2^M$ values. Each earlier
stopping index contributes a disjoint frozen copy of a state set of that
size, and there are at most $T+1$ stopping indices, including zero. The
stated count follows. The frozen pair retains exactly the information
allowed to determine the postponed output, and ignored fresh samples add
no information to it.
\end{proof}

\begin{lemma}[Conditional priors at finite events]\label{lem:conditional-priors}
Let $S\sim\sigma_d$ and let $H$ be a finite-valued random variable jointly
distributed with $S$. Every positive-probability value $h$ gives a
conditional signal density bounded by
$L_h=1/\Prob\{H=h\}$ relative to $\sigma_d$, and
\begin{equation}\label{eq:rare-event-entropy}
 \E\log L_H=H(H).
\end{equation}
\end{lemma}

\begin{proof}
Bayes' formula gives the Borel density
\[
 \frac{dP_{S\mid H=h}}{d\sigma_d}(s)
 =\frac{\Prob\{H=h\mid S=s\}}{\Prob\{H=h\}}\le L_h.
\]
Averaging $\log L_h$ gives \eqref{eq:rare-event-entropy}. In particular,
$H$ may be one padded state or a history of padded states; these have
different entropy costs.
\end{proof}

\begin{lemma}[Output capacity and the early clock bound]\label{lem:early-clock}
Consider a family of learners indexed by $d$, with uniform-prior success
at least a fixed $p_0>0$, and set
$\beta=\log(1/\epsilon)$ for $0<\epsilon\le1/10$.
At each dimension, the terminal output capacity satisfies
\begin{equation}\label{eq:output-capacity}
 p_0\le (T+1)2^M\epsilon^{d-1}.
\end{equation}
On the branch $T\le d\beta$, if $M=o(d^2)$, then
\begin{equation}\label{eq:early-clock}
 \beta=o(d),\qquad T=o(d^2),\qquad h_T=o(d^2).
\end{equation}
\end{lemma}

\begin{proof}
First fix a shared seed in the original conditional experiment.
There are at most $(T+1)2^M$ terminal state/index pairs. For a fixed pair,
write $\nu_{j,v}$ for its prescribed output law and
$r_{j,v}(s)$ for the conditional probability of reaching it given $S=s$.
Its contribution to success is
\[
 \int r_{j,v}(s)\int
 \ind_{\{\arccos\langle u,s\rangle\le\epsilon\}}
 \,\nu_{j,v}(du)\,\sigma_d(ds)
 \le\epsilon^{d-1},
\]
because $r_{j,v}\le1$ and each angular cap has mass at most
$\epsilon^{d-1}$ by Lemma~\ref{lem:spherical-ball}.
Summing gives \eqref{eq:output-capacity} for its
conditional success. The same numerical bound holds for almost every seed,
so averaging the original conditional success probabilities gives the
stated bound for the jointly measurable mixture.

On $T\le d\beta$, the inequalities $\beta\ge\log10$ and
$T+1\le2d e^\beta$ imply
\[
 (d-2)\beta\le M\log2+\log(2d)+\log(1/p_0).
\]
Dividing by $d$ proves $\beta=o(d)$ for a subquadratic memory sequence.
The branch then gives $T=o(d^2)$, and \eqref{eq:padded-entropy} gives
$h_T=o(d^2)$.
\end{proof}

\begin{lemma}[A raw-sample endpoint]\label{lem:raw-sample-bound}
Let $0<\epsilon\le1/10$, and let an estimator use all exact measurements
$(G,GS)$ from $b\le\lfloor d/2\rfloor$ independent Gaussian rows, together
with independent randomness, when $S\sim\sigma_d$. Its angular success
probability is at most
\begin{equation}\label{eq:raw-sample-bound}
 \frac{8}{d+2}+(4\epsilon)^{d-b-1}.
\end{equation}
The separate Gaussian-norm estimate gives the also valid bound
$90/d+(4\epsilon)^{d-b-1}$. In particular, uniform-prior success at least
$3/5$ requires $T>\lfloor d/2\rfloor$ for all sufficiently large $d$,
even if all raw data up to the horizon are retained.
\end{lemma}

\begin{proof}
For $d\le2$, both numerical bounds exceed one. We may therefore assume
$d\ge3$, when $b\le\lfloor d/2\rfloor\le d-2$.
Condition on all the data. The estimator and the signal are then
conditionally independent, and the conditional signal is the residual
sphere in Lemma~\ref{lem:residual-sphere}. On radius at least $1/2$,
the successful points have relative chordal cap mass at most
$(4\epsilon)^{d-b-1}$, by the argument in
Lemma~\ref{lem:residual-information}. Add the corresponding failure
probability from \eqref{eq:residual-chebyshev}. The zero-row case uses
the deterministic radius one.

If a learner stops by $T\le\lfloor d/2\rfloor$, grant it all of the first
$\lfloor d/2\rfloor$ rows and labels, including unused ones. Its output is
a kernel of those data and independent randomness. The displayed bound
tends to zero uniformly for $\epsilon\le1/10$, contradicting success
$3/5$. This proves the last assertion. The angular-cap version used in
the first replica route follows as well by replacing $4\epsilon$ by
the weaker $5\epsilon$.
\end{proof}
